%!TEX root = ../../thesis-trier.tex

Here, we show that for commutative $n$-state input automata
and an arbitrary regular constraint language,
the constrained synchronization problem is in \NL\ (\cref{thm:main_theorem_comm_input_NL})
and solvable in time $O(|\Sigma|^2n^{|\Sigma|+2})$ (\cref{thm:main_theorem_comm_input}).
For unary automata (which are commutative),
by \cref{constr_sync:thm:infinite,constr_sync:thm:finite},
the problem is \LOGSPACE-complete
when the constraint language is infinite and in $\AC^0$ otherwise.

\begin{theorem}
\label{thm:main_theorem_comm_input_NL}
 Fix a regular language $L \subseteq \Sigma^*$.
 Then, for a commutative input $n$-state automaton~$\mathcal A = (Q, \Sigma, \delta)$,
 the problem whether $\mathcal A$ admits a synchronizing
 word in $L$ is in \NL.
\end{theorem}
\begin{proof}
 Let $\Sigma = \{a_1, \ldots, a_k\}$ and $\automat{B}{P}{\Sigma}{\mu}{p_0}{F}$
 be a PDFA such that $L = L(\mathcal B)$.
 We give an algorithm that uses a memory cell to store a state from $\mathcal B$ initialized with~$p_0$
 and $k$ counters $c_1, \ldots, c_k$.
% The algorithm repeats the following passes for at most $|P| \cdot \binom{n}{2}$ steps, where 
% in the $i$-th step a letter $x_i \in \Sigma$ is computed in the following way: 
 The algorithm %iterates through $i \in \{1,2,\ldots,|P| \cdot \binom{n}{2}\}
 repeats at most $|P| \cdot \binom{n}{2}$ times the following steps
 to compute, in the $i$-th step, a letter $x_i \in \Sigma$:
 \begin{enumerate}
 \item  Guess the letter $x_i \in \Sigma$,
 update the state $p$ stored in the memory cell with $\mu(p, x_i)$
 and add one to the counter $c_j$ where $x_i = a_j$, $j \in \{1,\ldots,k\}$.
% Then it is $c_r = |x_1 \cdots x_i|_{a_r}$
% for all $r \in \{1,\ldots,k\}$.
 	
 \item \label{constrinput:algo:parttwo} If $\mu(p, x_i) \in F$,
 then iterate through the states $q \in Q$,
% and the counters $c_j$, $j \in \{1,\ldots,k\}$, compute
 compute the state
 $
 \applytrans{\delta}{q}{a_1^{c_1} \cdots a_k^{c_k}},
 $
 and check whether this always results
 in the same state. 
 If this is the case,
 then $a_1^{c_1}\cdots a_k^{c_k}$
 synchronizes $\mathcal A$.
 We have %$\psi(x_1 \cdots x_i) = \psi(a_1^{c_1}\cdots a_k^{c_k})$
 $|x_1 \cdots x_i|_{a_j} = c_j$ for each $j \in \{1,\ldots,k\}$
 and, since $\mathcal A$ is commutative and so its set of synchronizing
 words is commutative by 
 \crefitem{prelim:classes:comm:prop:set-sync-words}{prelim:classes:comm:prop:set-sync-words:1},
 $a_1^{c_1}\cdots a_k^{c_k}$
 synchronizes $\mathcal A$ if and only if $x_1 \cdots x_i$ synchronizes~$\mathcal A$.
 So, for $\mu(p, x_i) \in F$ and when $a_1^{c_1} \cdots a_k^{c_k}$
 synchronizes $\mathcal A$, we can abort the procedure and return \rYES.
 Otherwise, we continue with the next step or, for $i > |P| \cdot \binom{n}{2}$,
 abort the procedure and return \rNO. 
\end{enumerate}
 
 Part~\ref{constrinput:algo:parttwo} of the above algorithm
 can be realized by using a memory cell to compute $\applytrans{\delta}{q}{a_1^{c_1} \cdots a_k^{c_k}}$
 on-the-fly, an additional memory cell
 to store a resulting state $\applytrans{\delta}{q}{a_1^{c_1} \cdots a_k^{c_k}}$
 for comparison between different passes,
 a counter to iterate through the states $q \in Q$,
 a counter to iterate through $c_1, \ldots, c_k$
 and a counter to iterate from zero to $c_j$ for each $j \in \{1,\ldots,k\}$.
 By using a binary encoding all counters can
 be stored using logarithmic space.
 
	By~\cref{constr_sync:prop:pto_comm_rho_automata},
	a commutative automaton is a $\rho$-automaton.
	So, by \cref{constr_input:prop:bound_rho}, 
	if there exists a synchronizing word in~$L$,
	then there exists one of length at most~$|P|\cdot \binom{n}{2}$.
	Therefore, by nondeterminism, 
	the above algorithm finds it, which shows its correctness.
\end{proof}

	\begin{theorem}
		\label{thm:main_theorem_comm_input}
		Fix a regular language $L \subseteq \Sigma^*$.
		Then, for a commutative input automaton~$\mathcal A$
		with $n$ states,
		the problem whether $\mathcal A$ admits a synchronizing
		word in $L$ is solvable in time $O(|\Sigma|^2 n^{|\Sigma|+2})$.
	\end{theorem}
	\begin{proof}
		By Proposition~\ref{prop:aut_for_sync_words} we
		can compute in time $O(n^{|\Sigma|+2})$ an
		automaton~$\mathcal C$ of size $n^{|\Sigma|}$
		recognizing the set
		of synchronizing words of $\mathcal A$.
		Let $\automat{B}{P}{\Sigma}{\mu}{p_0}{F}$
		be an automaton such that $L = L(\mathcal B)$.
		By \cref{prelim:eqn:product-intersection}, the product automaton of $\mathcal C$ and $\mathcal B$
		accepts $L(\mathcal C) \cap L$.
		and can be computed in time $O(|P| |\Sigma|^2 n^{|\Sigma|})$.
		Then, checking that the resulting automaton
		for the intersection accepts a non-empty
		language can also be done in  
		time $O(|P||\Sigma|^2n^{|\Sigma|})$ 
		where $O(|P||\Sigma|^2n^{|\Sigma|}) = O(|\Sigma|^2n^{|\Sigma|})$ since $\mathcal B$ is fixed,
		as we simply have to check whether there exists a path from the start state to a final state
		which can be done by a breadth-first search algorithm~\cite{cormen01introduction}.
	\end{proof}

The degree of the polynomial measuring the running time depends on the alphabet size. But note that a fixed
constraint automaton also fixes the alphabet, hence this parameter is not allowed to vary in the input semiautomata.